% Abstract — softened claim per TSE reviewer expectations.
% Old: "eliminates non-deterministic test variance"
% New: "reduces non-deterministic test variance by two orders of magnitude in
%       an industrial Kubernetes setting"

\begin{abstract}
Pull-request preview environments on Kubernetes have become a popular
substrate for end-to-end testing, but the per-test database state they
expose is widely reported as a major source of test non-determinism. We
report on a controlled study of a 12k-LoC Kubernetes operator we
designed and deployed in production, which exposes a
\texttt{Preview} custom resource and supplies each preview with a
PostgreSQL database restored from a pre-warmed checkpoint instead of a
migration replay.

We evaluate five open-source applications spanning four runtime
ecosystems (Python/Flask, Go, Python/Django, TypeScript/Prisma,
Java/Spring) on a 3-node Azure Kubernetes Service cluster. Across
1{,}200+ test executions, we measure four research questions:
flakiness rate (RQ1), cross-PR concurrency invariance (RQ2), wall-clock
cycle time (RQ3), and bug-detection capacity preservation (RQ4).
We compare two configurations of the same operator that differ only in
the database setup step: \texttt{mode=restore} (the proposed
checkpoint-based mechanism) and \texttt{mode=migration} (the baseline
in widespread industrial use).

Our results show that the checkpoint-based isolation
\emph{reduces non-deterministic test variance by two orders of
magnitude} (failure-rate variance ${\sim}0$ vs $\approx 0.25$;
Cohen's $h$ at the theoretical maximum on 4 of 5 subjects; Cliff's
$\delta=-1$) and \emph{shortens the per-PR test cycle by a median of
2.5$\times$ on dynamically-typed stacks and up to 16$\times$ on
Spring-Boot applications} (Cohen's $d \in [2.4, 20.9]$; observed power
$\ge 0.99$). A two-one-sided-tests equivalence analysis suggests
isolation \emph{does not impair} bug-detection capacity within a 5
percentage-point margin where coverage permits the test. We release
the operator (Apache-2.0), the experimental harness, all 1{,}200+ raw
test outcomes, and a one-command Docker-based reproduction package.

We position these findings as a \emph{lower bound} on the achievable
speed-up and an \emph{upper bound} on residual flakiness in similar
ephemeral-environment architectures, and we discuss the construct,
internal, external, and conclusion-validity boundaries that frame the
generalization scope.
\end{abstract}

% Softened-title proposal:
% Old: "Checkpoint-based Database Isolation Eliminates Non-deterministic Test Variance in Kubernetes Preview Environments"
% New: "Checkpoint-based Database Isolation Reduces Non-deterministic Test Variance by Two Orders of Magnitude in an Industrial Kubernetes Setting: A 5-Stack Empirical Study"
